\section{Finite-step observability and nondegenerate refinements}
\label{app:observable}
The candidate has its own exact distribution at each time, which may
differ from the target. We use two numerical runs to bound how much
acceleration can change that exact distribution. Lemma~\ref{lem:weak}
controls one test average at a time. To control the full component of
acceleration that changes the distribution, we must also bound the
part that the chosen tests leave unmeasured.\AIassistance

\subsection{How the interval splits determine the comparison}
Adding more time points need not substantially change a numerical run.
For example, a split can leave one subinterval almost as long as the coarse
interval. The refinement energy below measures the reduction in squared
step sizes. A positive lower bound makes the two runs different enough to
reveal the leading error.

Split the $k$th coarse interval into positive fractions
$\theta_{k,1},\ldots,\theta_{k,m_k}$ summing to one. Define
\begin{equation}
 \beta_k=1-\sum_r\theta_{k,r}^2,\qquad \beta_k\ge\beta_0>0.
 \label{eq:refinement-energy}
\end{equation}
The error $E_r$ is measured at common coarse nodes; set $E=E_c+E_r$.
For any transported test with acceleration pairing $I(t)$, Lipschitz constant
$M$, and weak-remainder constant $C_w$, direct subtraction gives
\begin{equation}
 Z_j-Z_i=-\frac12\sum_{k=i}^{j-1}\beta_k h_k^2 I(t_k)+R_{ij},
 \qquad |R_{ij}|\le C_\times H\sum_{k=i}^{j-1}h_k^2,
 \quad C_\times=2C_w+M/2.
 \label{eq:balanced-contrast}
\end{equation}
Indeed the fine leading sum is $-\tfrac12\sum_{k,r}
\theta_{k,r}^2h_k^2 I(t_k+s_{k,r})$, with $0\le s_{k,r}\le h_k$.
Replacing each pairing by $I(t_k)$ costs at most $(M/2)\sum h_k^3$.
The two recursive remainders together cost at most $2C_wH\sum h_k^2$.
No numerical trajectory is restarted at an exact marginal. For equal
$m$-fold subdivision a smaller constant is available:
\begin{equation}
 \beta_0=1-1/m,\qquad
 C_\times=(1+m^{-2})C_w+\frac{m-1}{4m^2}M.
 \label{eq:equal-m-constant}
\end{equation}

\begin{lemma}[Explicit observation of a signed residual]
\label{lem:master-observer}
Suppose $|I|\le A_o$, $\operatorname{Lip}(I)\le M$, $|Z_k|\le DE$,
and \eqref{eq:balanced-contrast} holds. Assume $A_o,M>0$ and put
\[
 c_o=\min\{T/(2A_o),1/(2M)\},\qquad
 L_o=\max\{4/c_o,8C_\times/\beta_0\}.
\]
Then, on any coarse grid,
\begin{equation}
 \|I\|_\infty\le L_oH+
 \left(\frac{64DN E}{\beta_0c_o^2}\right)^{1/3}.
 \label{eq:observer-cube}
\end{equation}
On a uniform coarse grid this improves to
\begin{equation}
 \|I\|_\infty\le L_oH+
 \left(\frac{32DN E}{\beta_0Tc_o}\right)^{1/2}.
 \label{eq:observer-square}
\end{equation}
Zero values of $A_o$ or $M$ are handled directly or by increasing them to
any positive upper bounds.
\end{lemma}
\begin{proof}
The proof uses the fact that $I$ cannot change arbitrarily quickly.
If its magnitude is large at one time, it stays away from zero with the
same sign on a nearby interval. The difference between the two runs on
that interval then forces a lower bound on their target errors.

Let $a=\|I\|_\infty>0$ and choose a maximizing time. On its longer side
there is an interval of length
$\ell=\min\{T/2,a/(2M)\}\ge c_oa$ where $I$ has one sign and
$|I|\ge a/2$. If $H\le\ell/4$, its first and last coarse nodes are
separated by at least $\ell/2$. If also $H\le\beta_0a/(8C_\times)$,
\eqref{eq:balanced-contrast} gives
\[
 2DE\ge |Z_j-Z_i|\ge\frac{\beta_0a}{8}\sum_{k=i}^{j-1}h_k^2
 \ge\frac{\beta_0a\ell^2}{32N}.
\]
This proves the cubic bound. On a uniform grid the last squared-step sum
is at least $T\ell/(2N)$, proving the square-root bound. If either mesh
condition fails, $a\le L_oH$. Adding the two alternatives proves the
statements for every grid, including coarse grids where the bound is weak.
\end{proof}

For source tests with $\|\nabla\phi\|_\infty,
\operatorname{Lip}(\nabla\phi)\le1$, the constants in
Supplementary Section~\ref{app:uniform} and Lemma~\ref{lem:weak} can be used directly.
Taking the supremum over all such tests bounds the largest effect detected
by this test class at any time. This is a weak acceleration residual.
It is not yet an $L^2$ bound on the whole acceleration component.

\subsection{Bounding the full residual with a Gaussian polynomial expansion}
A finite collection of measurements can miss a nonzero acceleration
component. We therefore also bound the part that is not measured.
We express the transported acceleration as a sum of orthogonal Gaussian
polynomial components. A derivative estimate bounds the contribution of
the omitted higher-degree components, called the spectral tail.
This argument uses standard Gaussian initialization. Write
\[
 X_t=F_{0,t}(G),\quad J_t=DF_{0,t},\quad
 V_t=J_t^{-1}a_t\circ F_{0,t},\quad
 \varepsilon_t=\|\Pi_{\rho_t^u}a_t\|_{L^2(\rho_t^u)}.
\]
Here $J_t$ is the spatial Jacobian, which describes how the flow changes
when its initial position changes. The field $V_t$ expresses acceleration
in those initial coordinates. The projection $\Pi_{\rho_t^u}$ keeps the
component detected by gradients of test functions. Its norm $\varepsilon_t$
measures the part that can change the current distribution.
We use the constants $C_F,A,B,C_D,C_E,M_2,M_4$ from
\eqref{eq:base-constants}--\eqref{eq:flow-bounds}. Define
\begin{align}
 A_o&=AC_F^2\sqrt{M_2},&M&=BC_F^3\sqrt{M_4},\\
 c_o&=\min\{T/(2A_o),1/(2M)\},&
 V_1&=KTAC_F^5+C_D C_F^3. \label{eq:spectral-core}
\end{align}
Let $H_\alpha$ be normalized multivariate probabilists' Hermite polynomials,
so $\E H_\alpha(G)H_\eta(G)=\mathbf1_{\alpha=\eta}$. These polynomials are
orthogonal under the standard Gaussian law, so the mean product of two
different basis polynomials is zero. The multi-index $\alpha$ lists the
degree in each coordinate, and $|\alpha|$ is their total degree.
For $\ell=|\alpha|\ge1$ set $\phi_\alpha=H_\alpha/\sqrt \ell$.
In the polynomial envelopes below, $m=\ell$ denotes this total degree;
it is distinct from the subdivision factor in \eqref{eq:equal-m-constant}.
If
$\phi_\alpha(x)=\sum_b c_bx^b$, use the explicitly finite envelopes
\[
 P_{1,\alpha}=\sum_b|c_b||b|,\qquad
 P_{2,\alpha}=\sum_b|c_b||b|(|b|-1).
\]
Then $|\nabla\phi_\alpha(x)|\le P_{1,\alpha}(1+|x|)^{m-1}$ and
$\|D^2\phi_\alpha(x)\|\le P_{2,\alpha}(1+|x|)^{\max(m-2,0)}$.
Set $\mu_r=\E(1+|G|)^r$, including $\mu_0=1$. Every $\mu_r$ is explicit:
\[
 \mu_r=\sum_{j=0}^r\binom{r}{j}2^{j/2}
       \frac{\Gamma((d+j)/2)}{\Gamma(d/2)}.
\]
The next constants keep track of how each polynomial and its derivatives
grow after transport by the flow. They make the coefficient bounds
explicit. For each $\alpha$ define
\begin{align}
 G_\alpha&=P_{1,\alpha}C_F^\ell,&
 H_\alpha^*&=(P_{2,\alpha}+KTP_{1,\alpha}C_F)C_F^\ell,\\
 D_\alpha&=G_\alpha 2^{m-1}C_F^{m-1}\sqrt{\mu_{2m-2}},\\
 C_{0,\alpha}&=JC_F^3G_\alpha/6+KAC_FG_\alpha/2
             +TH_\alpha^*A^2C_F^{\ell+1}/8,\\
 C_{b,\alpha}&=H_\alpha^*A+G_\alpha C_D,\\
 C_{w,\alpha}&=\mu_{\ell+1}\{C_{0,\alpha}C_F^{\ell+1}
                   +\tfrac12 C_{b,\alpha}C_F^mC_ET\}.
 \label{eq:polynomial-constants}
\end{align}
Use \eqref{eq:balanced-contrast} or \eqref{eq:equal-m-constant} to define
$C_{\times,\alpha}$, and put
\begin{align}
 L_\alpha&=\max\{4/c_o,8C_{\times,\alpha}/\beta_0\},\\
 b_\alpha(E,H)&=L_\alpha H+
       \left(\frac{64D_\alpha NE}{\beta_0c_o^2}\right)^{1/3}.
 \label{eq:hermite-coefficient-bound}
\end{align}
On a uniform coarse grid replace the cube-root term by
$\{32D_\alpha NE/(\beta_0Tc_o)\}^{1/2}$.

\begin{theorem}[Explicit spectral residual observability]
\label{thm:spectral-proof}
For every $u\in\overline{\mathcal B}_K$, every balanced refinement,
and every integer $q\ge0$,
\begin{equation}
 \boxed{\quad
 \sup_{0\le t\le T}\varepsilon_t
 \le C_F\left\{
   \sum_{1\le|\alpha|\le q+1}b_\alpha(E,H)^2
       +\frac{dM_2 V_1^2}{q+1}\right\}^{1/2}.
 \quad}
 \label{eq:spectral-bound}
\end{equation}
The infimum over $q$ is a fully specified upper bound. The acceleration and
projection refer to the field's own exact law; exact target-path
preservation is not assumed. The same statement is valid if $E$ is replaced
by the maximum cross-grid $W_2$ distance, which may be smaller than $E_c+E_r$.
\end{theorem}
\begin{proof}
We first apply the signed observer to each polynomial coefficient.
We then bound the coefficients beyond the retained degree and carry
the resulting norm bound from source coordinates to the current law.
This second step is what turns finitely many coefficient estimates into
a bound for the full residual.

For $\phi=\phi_\alpha$, let $\phi_t=\phi\circ F_{t,0}$.
Its gradient and Hessian are bounded by
$G_\alpha(1+|x|)^{m-1}$ and $H_\alpha^*(1+|x|)^{m-1}$.
The second bound follows from the two terms in the second derivative of a
composition and $\|D^2F_{t,0}\|\le KTC_F^2$.
A Taylor expansion of $\phi_{t+h}$ at the exact endpoint, comparing it to
the Euler endpoint, gives a remainder at most
$C_{0,\alpha}h^3(1+|x|)^{\ell+1}$.
The three terms in \eqref{eq:polynomial-constants} bound three sources of
error. They are the change in particle acceleration, the transport of the
test gradient, and the Taylor remainder involving the Hessian. The Hessian
contains the test function's second spatial derivatives. For
example, along the segment between the two endpoints,
$1+|z|\le C_F(1+|x|)$ and
$|E-Y|\le A C_F h^2(1+|x|)/2$; its Hessian remainder is therefore at most
$H_\alpha^*A^2C_F^{\ell+1}h^4(1+|x|)^{\ell+1}/8$.

The pairing $\nabla\phi_t\cdot a_t$ has spatial derivative bounded by
$C_{b,\alpha}(1+|x|)^\ell$. Comparing recursive numerical states to exact
states and then summing over any subinterval gives exactly
$C_{w,\alpha}H\sum h_k^2$, using
$\E(1+|G|)^{\ell+1}=\mu_{\ell+1}$. Thus all recursive estimates continue to hold
for these unbounded polynomial tests; no target moment beyond the
Wasserstein comparisons is used.
For two numerical laws, the mean-value inequality and Cauchy--Schwarz give
\[
 |\E\phi_t(\widehat X)-\E\phi_t(\widetilde X)|
 \le D_\alpha W_2(\Law\widehat X,\Law\widetilde X).
\]
To check the constant, the squared gradient factor has moment at most
$G_\alpha^2 2^{2m-2}C_F^{2m-2}\mu_{2m-2}$ under any coupling, using only
the separate Gaussian-start marginal moment bounds.

The Hermite normalization gives $\|\nabla\phi_\alpha\|_{L^2(\gamma_d)}=1$.
The coefficient
$c_\alpha(t)=\E\nabla\phi_\alpha(G)\cdot V_t(G)$ therefore satisfies
$|c_\alpha|\le A_o$ and $\operatorname{Lip}(c_\alpha)\le M$, because
$|V_t(G)|\le AC_F^2(1+|G|)$ and
$\partial_tV_t=J_t^{-1}(j_t-Du_ta_t)\circ F_{0,t}$ has $L^2$ norm at
most $BC_F^3\sqrt{M_4}$. Lemma~\ref{lem:master-observer} yields
$\sup_t|c_\alpha(t)|\le b_\alpha$.

We must also account for components outside the retained finite set.
The gradients $\nabla H_\alpha/\sqrt{|\alpha|}$ form an orthonormal basis
of the Gaussian gradient subspace. Thus their coefficients account for
the full projected field, not only the measurements kept in the sum. Orthogonality follows
from $\partial_iH_\alpha=\sqrt{\alpha_i}H_{\alpha-e_i}$.
Completeness follows by expanding a smooth compact potential in Hermite
polynomials in the Gaussian Sobolev norm; integration by parts identifies
its derivative coefficients. Polynomial gradients themselves belong to
the closure of compact smooth gradients by Gaussian-tail cutoffs.
Gaussian differentiation also gives
$\sum_{i,b}|b|\,|\langle V_i,H_b\rangle|^2=\E\|DV\|_{\mathrm F}^2$ for
$V\in W^{1,2}(\gamma_d)$, first for polynomials and then by approximation in that Sobolev norm. Here the Gaussian Sobolev norm
controls both the square-integrable function and its square-integrable first
derivatives under the Gaussian law.
Since projection onto gradients preserves each total Hermite degree,
\[
 \|\Pi_{\gamma_d}V_t\|_2^2
 \le\sum_{1\le|\alpha|\le q+1}|c_\alpha(t)|^2
       +\frac{\E\|DV_t\|_{\mathrm F}^2}{q+1}.
\]
The chain rule, $D(J^{-1})=-J^{-1}(DJ)J^{-1}$, and the flow estimates give
$\|DV_t(x)\|_{\rm op}\le (KTAC_F^5+C_D C_F^3)(1+|x|)$, hence the numerator
is at most $dM_2V_1^2$.
Finally, for any physical test $\psi$ and $\phi=\psi\circ F_{0,t}$,
\[
 \langle\nabla\psi,a_t\rangle_{\rho_t^u}
 =\langle\nabla\phi,V_t\rangle_{\gamma_d},\qquad
 \|\nabla\phi\|_2\le C_F\|\nabla\psi\|_{L^2(\rho_t^u)}.
\]
Taking the dual supremum proves
$\varepsilon_t\le C_F\|\Pi_{\gamma_d}V_t\|_2$ and the theorem.
All derivative estimates are valid almost everywhere for closure fields; the weak
expansion passes from smooth approximants with the stated uniform
remainders.
\end{proof}

The cutoff degree $q$ sets how many coefficients enter the finite sum.
For fixed $q$, all constants are finite and explicit. They can nevertheless
grow rapidly with $q$, $K$, and dimension, so the bound is not a practical
dimension-free certificate.

The order of choices explains what the bound does establish. First choose
$q$ large enough to make the tail small. Then make $H$ and $NE$ small
enough to control the retained coefficients. Consequently, $H\to0$ and
$NE\to0$ force the \emph{full} projected-acceleration norm to zero,
uniformly in time. This conclusion does not assume that a finite dictionary
alone detects every component. It also gives no power-law estimate of
distance to the set of exact generators.

\paragraph{Converting a residual lower bound into an error lower bound.}
Suppose a separate argument proves that the projected acceleration cannot
be arbitrarily small. The upper bound just obtained then prevents the
two-grid error from being arbitrarily small. The construction below gives
explicit constants for this implication.

Fix a proved floor $\sup_t\varepsilon_t\ge\delta>0$. Choose an integer $q$
with $q+1\ge4C_F^2dM_2V_1^2/\delta^2$. Write
$b_\alpha=L_\alpha H+\Lambda_\alpha(NE)^{1/p}$, where $p=3$ for arbitrary
balanced grids and $p=2$ for uniform grids, and let
$L_q=(\sum L_\alpha^2)^{1/2}$ and
$\Lambda_q=(\sum\Lambda_\alpha^2)^{1/2}$ over the retained indices. Then
\[
 H\le H_\delta:=\frac{\delta}{4C_FL_q}
 \quad\Longrightarrow\quad
 E\ge\frac{c_\delta}{N},\qquad
 c_\delta=\left(\frac{\delta}{4C_F\Lambda_q}\right)^p.
\]
Indeed, a smaller error makes the retained coefficient norm at most
$\delta/(2C_F)$ and the tail at most $\delta^2/(4C_F^2)$, contradicting the
floor via \eqref{eq:spectral-bound}. All constants are explicit. A proved positive acceleration pairing
can supply such a floor; a target-feature discrepancy alone cannot.
The constants and number of retained coefficients may be impractical.

\subsection{Why bounded tests alone do not give the full norm}
A test average can be small because positive and negative contributions
cancel. The squared norm does not have this cancellation. A rapidly
oscillating acceleration can therefore have very small smooth-test
averages while its norm is larger. The example below shows why no common
linear bound compares these quantities under the stated derivative bound.

For $d=1$, $u_n(t,x)=t\sin(nx)/n^2$ belongs to $\mathcal B_1$ for
$0\le t\le1$, and $a_0=\sin(nx)/n^2$. Its full source residual is
\[
 \|a_0\|_2=\frac{\sqrt{(1-e^{-2n^2})/2}}{n^2}.
\]
In contrast, Gaussian integration by parts gives, for
$|\phi'|,|\phi''|\le1$,
\[
 |\E\phi'(G)a_0(G)|\le (1+\E|G|)/n^3.
\]
The flow of $u_n$ satisfies $F_{0,t}(x)=x+O(n^{-2})$ and
$DF_{0,t}(x)=1+O(n^{-1})$ uniformly for $0\le t\le1$, so the transported
pairings are $O(n^{-3})$ at every time.
Hence no common linear bound of the full norm by this weak residual holds,
even within the stated derivative bound. A finite fixed dictionary can
also miss an entire nonzero orthogonal component. The spectral tail in
\eqref{eq:spectral-bound}, rather than an unproved norm equivalence, repairs
that logical gap.

\section{Refinement and Runge--Kutta versions of the characterization}
\label{app:scheme-extension}
The characterization depends on separating the leading errors of the two
numerical runs. Balanced interval subdivision supplies that separation for
Euler. A suitable pair of Runge--Kutta methods can supply it on one common
grid. We state the conditions for each comparison and then give examples
showing why weaker conditions do not suffice.\AIassistance

With \eqref{eq:refinement-energy}, the necessity proof of
Theorem~\ref{thm:characterization} applies with the same steps after replacing its
contrast by \eqref{eq:balanced-contrast}. A nonzero limiting pairing
retains one sign on a fixed interval of length $\ell$, and $\sum h_k^2\ge\ell^2/(4N)$ there.
The constant $\beta_0$ must be uniform in the sequence. Sufficiency is
unchanged since the fine maximum step is at most $H$.
Thus all equivalences and the uniform $c/N$ conclusion hold for every
fixed integer refinement factor and every uniformly balanced, possibly
interval-dependent subdivision. Fine step counts may differ from $2N$;
the reference step count is the coarse number of steps.

\subsection{Pairs of Runge--Kutta methods with distinct leading defects}
A Runge--Kutta update evaluates the velocity at several intermediate stage
states before forming the next state. Its coefficients specify those states
and the weighted output. The leading defect is the first local error term
relative to the exact update. The scalar coefficient below determines that
term under the stated consistency condition.

Consider a fixed explicit Runge--Kutta method with coefficients
$\mathsf A_{ij}$, $\mathsf b_i$, $c_i=\sum_j \mathsf A_{ij}\in[0,1]$, and $\sum_i \mathsf b_i=1$.
All coefficients and the number of stages are independent of $N$.
Write
\[
 \kappa=\tfrac12-\sum_i \mathsf b_ic_i.
\]
The internal consistency $c=\mathsf A\mathbf1$ is essential when the
velocity depends explicitly on time. It makes the two second-order terms
involving time and space derivatives have the same coefficient.
Taylor expansion of the stages gives
\begin{equation}
 \Psi_h^u(t,x)=x+hu(t,x)+h^2(\tfrac12-\kappa)a_t(x)
                 +O(h^3(1+|x|)^2).
 \label{eq:rk-leading-defect}
\end{equation}
The remainder bound holds for every field in $\mathcal B_K$ with the same
constants. Applying the estimates stage by stage gives common linear-growth
and Lipschitz bounds $1+Ch$ for $h\le H_0$. The constants depend only on
$K,T$ and the fixed tableau. The $C^2$ derivative bound bounds each
Taylor remainder. Gaussian fourth moments then bound its $L^2$ norm.

\begin{theorem}[Leading-defect separation for a solver pair]
\label{thm:rk-pair-proof}
Run two such methods on the same deterministic $N$-step grid with the same
field, and let their constants satisfy $\kappa_1\ne\kappa_2$.
Replace $E_c+E_f$ in Theorem~\ref{thm:characterization} by the sum of
these methods' maximum-node target $W_2$ errors. All three equivalences
and the uniform $c/N$ alternative remain valid. The common derivative bound and
target path are unchanged; the fields and grids may depend on $N$.
\end{theorem}
\begin{proof}
For a transported bounded-gradient test, \eqref{eq:rk-leading-defect}
and the same recursive state comparison as Lemma~\ref{lem:weak} give
\[
 D_j^{(l)}-D_i^{(l)}=-\kappa_l\sum_{k=i}^{j-1}h_k^2 I(t_k)
             +O\!\left(H\sum_{k=i}^{j-1}h_k^2\right),\quad l=1,2.
\]
Subtracting cancels the model law and retains coefficient
$\kappa_1-\kappa_2\ne0$. All constants are common to the derivative class.
The constants can depend on the full fixed tableaus.
The strong first-order bound \eqref{eq:strong} and the node-to-path bound
\eqref{eq:bias-general} hold for each method with such constants, for
$H\le H_0$.
Strong first-order convergence on arbitrary refining grids gives an exact
limiting generator whenever the two target errors vanish. A nonzero acceleration
pairing in that generator persists on an interval; the signed contrast and
$\sum h_k^2\ge\ell^2/(4N)$ contradict superlinear two-method accuracy,
exactly as in Supplementary Section~\ref{app:characterization}.

Conversely, suppose an order-elevating generator belongs to the closure.
We change how numerical and exact particles are paired without changing
their distributions. For method $l$, move the exact particles along the
acceleration field for auxiliary time $\kappa_l h^2$, while physical time
stays fixed. This auxiliary time may be negative. Linear growth makes the
acceleration flow complete in both directions. Inserting the recoupled state into
\eqref{eq:rk-leading-defect} gives the exact trajectory through order two,
with local $W_2$ defect $O(h^3)$; the local Lipschitz and moment estimates
are the same as the Euler recoupling proof up to tableau constants.
Method stability propagates the error to $O(H^2)$. Smooth approximants to
a generator in the closure approximate both finite numerical runs simultaneously.
The contradiction proof of a common lower bound is then unchanged.
\end{proof}

Euler has $\kappa=1/2$ and Heun has $\kappa=0$, so Euler--Heun on the same
grid is a concrete instance, using $N$ and $2N$ field evaluations.
Two second-order Runge--Kutta methods both have $\kappa=0$ and do
\emph{not} satisfy this theorem: ordinary second-order accuracy already
allows superlinear two-grid accuracy for any sufficiently regular exact generator.
The theorem is a class of solver-pair converses, not a claim that every
integrator has the Euler obstruction.

Distinct leading defects are needed even when both methods have first-order
accuracy.
Take an explicit three-stage tableau with $c=(0,1/2,1)$, $\mathsf A_{21}=1/2$,
$\mathsf A_{31}=1$, other entries zero, and
$\mathsf b=(1+\eta,-2\eta,\eta)$ for fixed $\eta\ne0$.
It is consistent, has $\mathsf b^\top c=0$, and is distinct from Euler, but both
have $\kappa=1/2$. Their local difference is $O(h^3)$.
On the scalar Gaussian secant family Euler is target-exact; this distinct
method has $O(h^2)$ target error on the same uniform grid even though no
order-elevating generator exists. Thus distinct method names alone are not
sufficient either. Negative weights in this counterexample are permitted
by the stated class; no strong-stability-preserving property is claimed.

\paragraph{Degenerate refinement is a genuine counterexample.}
For the Gaussian secant field in \eqref{eq:secant}, split each step $h=1/N$
into $h/N$ and $h(1-1/N)$. There are $2N$ fine steps, but
$\beta_k=2N^{-1}(1-N^{-1})$. A direct second-order expansion of the two
updates shows their map differs from the coarse Euler update by
$O(h^3)(1+|x|)$, uniformly in $k$. Consequently both target errors are
$O(N^{-2})$ although no order-elevating generator exists in dimension one.
Likewise, a field fitted to $N$ uniform steps has $O(N^{-2})$ error on
$N+1$ uniform steps: their leading step sizes differ by only $O(N^{-2})$.
A statement about two \emph{distinct counts} without a nondegenerate
contrast is therefore false.

\paragraph{An additional assumption for parameter distance to a generator.}
The residual measures an effect of acceleration, not a distance between
parameter vectors. A small residual need not place the model parameters
close to those of an exact generator. That further conclusion requires
local identifiability in a finite-dimensional exact-path family.
Here identifiability means that every parameter direction changes the
selected observations by a quantitatively nonzero amount. Let $\theta\mapsto
u_\theta$ be such a family with the common derivative bound, let $R(\theta)$ collect
finitely many transported Hermite coefficients at fixed times, and suppose
$R(\theta_*)=0$ for an order-elevating generator,
$s_{\min}(DR(\theta_*))\ge\sigma>0$, and $DR$ is $L_R$-Lipschitz on a
parameter ball. For $\|\theta-\theta_*\|\le\sigma/L_R$, Taylor's theorem gives
\[
 \|\theta-\theta_*\|\le 2\|R(\theta)\|/\sigma.
\]
The observer bounds control each coordinate of $R$. The proof is simply
$\|R(\theta)\|\ge\sigma\|\theta-\theta_*\|-L_R\|\theta-\theta_*\|^2/2$.
This is a conditional inverse-function estimate, not an assertion that an
arbitrary neural parametrization is identifiable, that a nearby order-elevating generator exists,
or that the selected finite tests characterize all order-elevating generators globally.

\paragraph{Why first-derivative bounds alone do not suffice.}
The uniform bounds must control changes in time as well as changes in space.
Otherwise, a field can oscillate between the evaluated times and fit both
numerical grids without satisfying the required acceleration condition.
On the scalar nonaffine Gaussian path let $a_h(t)=[g(t+h)-g(t)]/[hg(t)]$,
with $g$ smoothly extended to a neighborhood of the time strip. Define
\begin{align}
 b_h(t)&=\frac{a_h(t-h/2)}{1+(h/2)a_h(t-h/2)},\\
 c_h(t)&=a_h(t)+\{b_h(t)-a_h(t)\}\sin^2(\pi t/h).
 \label{eq:temporal-alias-field}
\end{align}
For small $h$, the field $u_h(t,x)=c_h(t)x$ has common bounds on
$|u|/(1+|x|)$, $|\partial_tu|/(1+|x|)$, $D_xu$,
$D_x\partial_tu$, and $D_x^2u=0$. Indeed $b_h-a_h=O(h)$ in every fixed
coefficient derivative norm, so differentiation once of the oscillatory
factor remains bounded. At a coarse node $kh$, $c_h=a_h$ and the coarse
map is exactly the target scale ratio. At its midpoint,
$c_h(kh+h/2)=a_h(kh)/(1+ha_h(kh)/2)$, and the two half-step multipliers
multiply to $1+ha_h(kh)$. Both runs are therefore target-exact at \emph{every
common coarse node}. Nevertheless the path has no order-elevating generator in
dimension one. The second time derivative of $c_h$ is generally
$\Theta(h^{-1})$, violating the common derivative bound precisely where compactness
of the first time derivative is needed. This disproves a version of the
characterization that keeps only uniform first-derivative bounds. It does
not rule out replacing the second-derivative bound by an appropriate
uniform modulus of continuity.
